\section{Metodología de evaluación}\label{sec:dar-metodologia}

La metodología de evaluación empleada para la selección de la tecnología de servicio de directorio fue el \ac{DAR} de \ac{CMMI} \citep{CMMIInstitute2010}. Esta metodología permitió evaluar de manera estructurada las alternativas tecnológicas identificadas, tomando como base los requisitos establecidos para el servicio de directorio. Para ello, se definieron criterios de evaluación derivados de dichos requisitos y complementados con aspectos relacionados con la gestión de identidades y la seguridad de la información, considerando como referencia las normas ISO/IEC 24760 e ISO/IEC 27002.

La evaluación consideró aspectos relacionados con la autenticación, la gestión de cuentas y usuarios, la seguridad de las comunicaciones, la compatibilidad con la infraestructura existente, la administración de accesos y la interoperabilidad, entre otros. Cada criterio fue ponderado de acuerdo con su importancia para el proyecto y las alternativas fueron valoradas según su nivel de cumplimiento. A partir de estas valoraciones se obtuvo una puntuación ponderada para cada alternativa, permitiendo establecer una comparación objetiva y fundamentar la selección de la tecnología con mayor correspondencia con los requisitos definidos.
